\section{Local energy bounds with retained excitations}\label{en:section}\label{sec:energy}

A scalar energy asymptotic does not by itself determine the one-body
momentum distribution.  We need a lower bound that retains the positive
quadratic excitation operator as well as the leading energy.  This
section proves that bound on boxes of side $E$, with $E=B$ or $R$;
Lemma~\ref{en:sharp} is the precise low-excitation statement.
Section~\ref{en:consequences} will compare the summed lower bound with
the ground-state energy upper bound and extract count fluctuations and
local one-body covariances.  The remaining occupation of wavelengths
comparable with the whole torus requires the spatial argument later in
the paper.

Throughout this section, global expectations are in a normalized bosonic
ground vector, with no sign or reality restriction.  Local inequalities
also apply to unnormalized box restrictions with spectator variables.
For a cell $A$, write $n_A=\sum_{i=1}^N\1_A(x_i)$ and
$q_A=n_A/(D|A|)$.  Constants are uniform for $\alpha$ in the fixed compact
window of Section~\ref{set:section}; they may depend on the fixed
exponents and, when indicated, on $j$.

\subsection{The local energy and the relevant scales}

Put $c_{\rm L}=128/(15\sqrt\pi)$ and, with $r$ held fixed, define
\begin{equation}\label{en:energy-density}
 e_r(x)=4\pi r x^2+4\pi c_{\rm L}r^{5/2}x^{5/2}.
\end{equation}
The hard-sphere upper bound of \cite[Theorem~1.1]{BBCOS}, after rescaling,
gives
\begin{equation}\label{en:upper}
 E_{N,K,r}\le V\{e_r(D)+O(D^{-c_*})\}
\end{equation}
for an absolute $c_*>0$, in the order of limits just specified.  That
theorem uses periodic boxes and a fixed hard-core scattering length.
To obtain this upper bound along a density-convergent sequence, one may
fix a slightly higher physical density $\rho'>\rho$ and choose, in the
same physical box, an integer $M\ge N$ with $M/L^3\to\rho'$.
Slicing an $M$-particle trial state and then averaging over its omitted
coordinates gives $E_{N,L,a}\le E_{M,L,a}$.  The thermodynamic upper
bound at $\rho'$ therefore bounds the limsup at $\rho$.  Its explicit
right side is continuous in $\rho'$, so one may let $\rho'\downarrow\rho$.
Equivalently, one may use the sequence-independent thermodynamic energy
limit.  No uniform finite-volume rate is used here.

We use a generic box side $E\asymp Q=D^\theta$, where $\theta=u$ or $w$,
and write $v_b=E^3$.  Neumann restriction to the boxes lowers the energy.
More explicitly, decompose configuration space into the sectors assigning
each label to a box, restrict each coordinate integral to its assigned
box, and discard exclusions between different boxes.  This operation
takes no derivative of a sector indicator.  Variables assigned elsewhere
remain spectators, so every local form inequality below may be integrated
in those variables.  The affine comparison energy for a box containing
$n$ particles is
\begin{equation}\label{en:tangent}
 t_b(n)=v_b e_r(D)+e_r'(D)(n-Dv_b).
\end{equation}
Its sum over all boxes is $Ve_r(D)$ in every assignment sector.

Use \cite[Proposition~2.1]{FJGMOT} to choose a nonnegative radial bounded
potential $v$ supported in the hard core.  Its scattering length $a'$ and
its scattering functions $\varphi$, $\omega=1-\varphi$, $g=v\varphi$
satisfy
\begin{equation}\label{en:replacement}
 0\le r-a'\le CrD^{-1}Q^{-1},\quad
 \|v\|_1\le CQ,\quad
 \|v\|_\infty\le D^{C},\quad
 \widehat g(0)=8\pi a',\quad
 \|g\|_1+\|g\omega\|_1\le C/D.
\end{equation}
The replacement also satisfies the radial comparison condition (2.6)
in that proposition.  The Neumann operator with this potential is denoted
by $H_b$.  Its expectation on the original hard-core restriction equals
the restricted kinetic energy, since $v$ vanishes outside the excluded
region.  It can subsequently be applied to excitation-localized vectors,
which need not satisfy the hard-core condition.

The fixed factors in the box lengths matter for the particle-number
hypothesis of \cite{FJGMOT}.  Use the choices
\eqref{set:scale-factors}--\eqref{set:nested-lengths}, so that
$E=b_\theta D^\theta(1+o(1))$ with
$b_\theta=\alpha_0^{-(1+3\theta)/2}$, in the specified order of limits.
Define
\[
 \rho_{\rm aux}
   =E^{-2/(1+\theta)}(a')^{-(1+3\theta)/(1+\theta)},\qquad
 K_\ell=(\rho_{\rm aux}(a')^3)^{-\theta/2}.
\]
Then the source length relation
$E=K_\ell/(\rho_{\rm aux}a')^{1/2}$ holds exactly.  Moreover,
\[
 \frac{\rho_{\rm aux}}D
 =\left(\frac{\alpha_0}{\alpha}\right)^{(1+3\theta)/(1+\theta)}
          (1+o(1))\in[1/3,3]
\]
for sufficiently small $\theta$ and sufficiently large $D$, uniformly
in the fixed $\alpha$ window.  The change from $r$ to $a'$ only affects
the $o(1)$ factor; the replacement bounds are chosen with fixed slack.
Consequently every
\begin{equation}\label{en:middle-sector}
 Dv_b/2\le n\le6Dv_b
\end{equation}
lies in the particle-number range $n\le20\rho_{\rm aux}v_b$ of that
paper.  Set
\begin{equation}\label{en:parameters}
 K_H=K_\ell^{100},\qquad
 M\asymp Dv_bQ^{-400}.
\end{equation}
Rounding $M$ to an integer changes no estimate.  For sufficiently small
$\theta$, the required inequalities are
\begin{gather}
 K_H\gg K_\ell^4,\qquad
 K_\ell K_H^3\ll(\rho_{\rm aux}(a')^3)^{-1/2},\qquad
 K_\ell^{5/4}K_H^2\ll(\rho_{\rm aux}(a')^3)^{-1/2},\label{en:parameter-conditions}\\
 1\ll M\ll E/a',\qquad
 M\ll \rho_{\rm aux}v_bK_H^{-3}K_\ell^{-17/4},\qquad
 K_H\ll\sqrt{E/a'}.\nonumber
\end{gather}
For example, the ratio in the second upper bound for $M$ is
$O(Q^{-95.75})$.  All powers containing $D$ rather than $Q$ can be made
favorable by decreasing $\theta$.  We shall also require
$u\ll c_*$ and several inequalities of this form below; they form a
finite list.

Let $u_p$ be the normalized cosine basis of the Neumann cube, indexed by
$p\in(\pi/E)\mathbb N_0^3$.  Thus
\[
 u_p(x)=v_b^{-1/2}\prod_{i=1}^3 c_{p_i}\cos(p_ix_i),\qquad
 c_0=1,\quad c_s=\sqrt2\ (s\ne0).
\]
Write $a_p,a_p^*$ for the associated annihilation and creation operators,
$n_+=\sum_{p\ne0}a_p^*a_p$, and
\begin{equation}\label{en:bands}
 \mathcal P_L=\{p:0<|p|\le K_H/E\},\quad
 n_L=\sum_{p\in\mathcal P_L}a_p^*a_p,\quad n_H=n_+-n_L.
\end{equation}
All estimates on unnormalized vectors include their squared norm in
every scalar error.

\subsection{A rough bound retaining kinetic energy and count fluctuations}
\label{en:rough-section}

The sharp box estimate will be used only after restricting the number of
low-frequency excitations.  We first establish a less accurate bound that
pays for this restriction and controls its complementary part.  The point
of the construction is to use most of the kinetic energy only near another
particle, leaving enough energy to control the low Neumann modes.  We use
Dyson's radial replacement and the small-box energy method of
Lieb and Yngvason~\cite{Dyson,LY}, keeping track of the unused kinetic
energy explicitly.

Throughout this subsection let \(\mathcal C=[0,E]^3\),
\(v_b=E^3\), and \(E\asymp Q=D^\theta\).  The exponent \(\theta>0\)
will be either of the two box exponents.  We assume that it is sufficiently
small, with bounds specified below.  Let \(r=\alpha/D\), where \(\alpha\)
ranges over the fixed compact window of positive values.  The nonnegative
radial replacement potential, denoted here by \(v\), is supported in the
ball of radius \(r\).  Its scattering length \(a'\) obeys
\[
  0\le r-a'\le CrD^{-1}Q^{-1},\qquad \|v\|_1\le CQ.
\]
Write
\[
 H_n=\sum_{i=1}^n(-\Delta_{\mathcal C,i}^{\mathrm N})
       +\sum_{1\le i<j\le n}v(x_i-x_j).
\]
All inequalities below are quadratic-form inequalities.  They hold with
additional Hilbert-space-valued spectator variables.  Let \(A\) run over
the cubes of side \(d\) in a fixed subdivision of \(\mathcal C\), and
put \(n_A=\sum_i\1_A(x_i)\).  The side \(d\) is fixed while
\(D\to\infty\).

\begin{lemma}[Rough box bound]\label{en:rough-bound}\label{en:rough}
There are constants \(c,C>0\) and an exponent \(c_{\mathrm r}>0\), independent
of the sufficiently small box exponent \(\theta\), such that, for
\(0\le n\le6Dv_b\),
\begin{equation}\label{en:rough-ineq}
 \begin{split}
 H_n\ge {}&\frac{4\pi r n^2}{v_b}-Cv_bD^{1-c_{\mathrm r}}
       +c\,\dd\Gamma\bigl(\min\{p^2,1\}\bigr)\\
 &\quad+cD\sum_A |A|\,
   \mathfrak h\left(\frac{n_A}{D|A|}-\frac{n}{Dv_b}\right),
 \qquad \mathfrak h(t)=\frac{t^2}{1+|t|}.
 \end{split}
\end{equation}
Here \(p\) denotes the Neumann frequency and \(\dd\Gamma\) sums
the indicated one-particle operator over all particles.  One may take
\(c_{\mathrm r}=10^{-3}\).
\end{lemma}

\begin{proof}
We give separately the short-distance energy estimate and the recovery of
the kinetic term.  Choose equal subdivisions of every \(d\)-cube with
side \(\ell\asymp D^{-0.31}\), and set
\[
 s\asymp D^{-0.36},\qquad \eps_D=D^{-0.002}.
\]
Rounding the number of subdivisions changes only fixed constants.  For
large \(D\), the support radius \(r\) is smaller than \(s/2\).
Restrict configuration space to sectors specifying which small cube
contains each labeled particle.  Restriction of an \(H^1\) function to
these sectors incurs no gradient cost; their faces carry Neumann form
conditions.  Discard interactions between different small cubes.

\emph{The radial replacement.}
Fix one particle coordinate and all the others in its small cube.  Partition
that cube into Voronoi regions of the other particles, clipping at the
cube boundary.  Each region is star-shaped about its center.  On a ray
from a center, the scattering variational principle gives, for \(R\ge r\),
\begin{equation}\label{en:radial-scattering}
 \int_0^R\left((1-2\eps_D)|f'(t)|^2
           +\tfrac12v(t)|f(t)|^2\right)t^2\,\dd t
       \ge (1-2\eps_D)a'|f(R)|^2.
\end{equation}
For completeness, with kinetic coefficient one the minimizer at fixed
endpoint is the regular scattering solution divided by its endpoint
value.  Its minimum is
\(a_v(1-a_v/R)^{-1}|f(R)|^2\), where \(a_v\) is the scattering
length of the potential in the integrand.  This follows by integrating
the radial scattering equation against the minimizer; the boundary term
is \(R^2\overline{f(R)}f'(R)\).  Completing the quadratic form around
the minimizer proves the same lower bound for every ray function.
Factoring out \(1-2\eps_D\) replaces \(v\) by
\(v/(1-2\eps_D)\), whose scattering length is at least \(a'\),
by the variational definition and monotonicity.  Dropping the factor
\((1-a_v/R)^{-1}\ge1\) proves
\eqref{en:radial-scattering}.  This argument applies first to smooth
functions and then by form approximation.

Let \(U\) be constant on \(s/2\le |x|\le s\), zero elsewhere,
and normalized by \(\int_{\R^3}U=4\pi\); thus
\(\int_0^\infty U(t)t^2\,\dd t=1\).
On a ray ending at distance \(T\), average
\eqref{en:radial-scattering} over radii \(R\le T\) with weight
\(U(R)R^2\,\dd R\).  The total weight is at most one, so only
the energy at distances at most \(s\) is used.  Integrating the rays
proves the lower bound
\((1-2\eps_D)a'U(d_{\mathrm{nn}})\), where
\(d_{\mathrm{nn}}\) is the nearest-neighbor distance within the
small cube.  If there is no other particle this expression is zero.
Half of every pair interaction is assigned to each endpoint, so summing
over the coordinates uses no interaction twice.

More explicitly, let \(S_i\) be the set of positions of coordinate
\(x_i\) within distance \(s\) of another particle in its small cube.
The preceding argument leaves the form
\begin{equation}\label{en:kinetic-reserve}
 \mathcal T_{\mathrm{rem}}[\psi]
 =\sum_i\int\left[\eps_D+(1-2\eps_D)\1_{S_i^c}\right]
                    |\nabla_i\psi|^2.
\end{equation}
There is also a separate reserve \(\eps_D\sum_i\|\nabla_i\psi\|^2\)
available for the next step.  Indeed, these two reserves and the
coefficient \(1-2\eps_D\) used inside \(S_i\) sum to the original
kinetic coefficient at every point.

\emph{The energy in a small cube.}
The cases \(k=0,1\) have zero scalar lower bound and are immediate.
For \(k\ge2\) particles in a cube of side \(\ell\), put
\[
 W_k=(1-2\eps_D)a'\sum_{i=1}^kU(d_{\mathrm{nn},i}).
\]
The first nonzero eigenvalue of the reserved kinetic operator is
\(\Delta_\ell=\eps_D\pi^2\ell^{-2}\).
Its normalized constant vector will be denoted by \(\Omega_k\).
For any fixed \(C_0\), uniformly for
\(k\le C_0D\ell^3\),
\begin{equation}\label{en:smallcube-mean}
 \langle W_k\rangle_{\Omega_k}
 \ge\frac{4\pi(1-2\eps_D)a'k(k-1)}{\ell^3}
       \left(1-C\frac{s}{\ell}
                 -C\frac{ks^3}{\ell^3}\right).
\end{equation}
To see this, keep the first coordinate a distance at least \(s\) from
the cube boundary.  For each possible nearest neighbor integrate that
neighbor over the shell of radii \(s/2,s\).  The other \(k-2\)
points lie outside the corresponding ball with probability at least
\(1-Cks^3/\ell^3\), by the union bound.  These events for different
nearest neighbors are disjoint up to null sets.  The removed boundary
layer occupies at most \(Cs/\ell\) of the cube volume.  Summing
over the first coordinate proves \eqref{en:smallcube-mean}.

The multiplication operator satisfies
\[
 0\le W_k\le Cka's^{-3},\qquad
 \langle W_k^2\rangle_{\Omega_k}
       \le Cka's^{-3}\langle W_k\rangle_{\Omega_k}.
\]
The scalar mean in \eqref{en:smallcube-mean} is less than
\(\Delta_\ell/2\) for large \(D\).  Temple's inequality therefore
gives
\begin{equation}\label{en:smallcube-temple}
 \inf\operatorname{spec}\left(
    \eps_D\sum_{i=1}^k(-\Delta_i^{\mathrm N})+W_k\right)
 \ge \langle W_k\rangle_{\Omega_k}
   -\frac{\langle W_k^2\rangle_{\Omega_k}}
          {\Delta_\ell-\langle W_k\rangle_{\Omega_k}}.
\end{equation}
This version of Temple's bound follows by decomposing the constant
trial vector into the ground eigenspace and its orthogonal complement;
the latter spectrum is at least \(\Delta_\ell\), by the min--max
principle and \(W_k\ge0\), and the variance bounds its spectral weight.
The powers used here are
\[
 D\ell^3\asymp D^{0.07},\quad
 s/\ell\asymp D^{-0.05},\quad
 Ds^3\asymp D^{-0.08},\quad
 ka's^{-3}\le CD^{0.15},\quad
 \Delta_\ell\asymp D^{0.618}.
\]
Consequently there is \(\delta_D\le CD^{-0.002}\) such that
the energy available in each small cube is bounded below by
\begin{equation}\label{en:smallcube-quadratic}
   \frac{4\pi r}{\ell^3}(1-\delta_D)k(k-1),
       \qquad 0\le k\le C_0D\ell^3.
\end{equation}
The replacement of \(a'\) by \(r\) is included in \(\delta_D\).

For larger \(k\), divide its labels into groups of sizes between
\(m/2\) and \(m\), where \(m=\lfloor C_0D\ell^3\rfloor\).
Such a partition exists for \(k>m\), by taking
\(\lceil k/m\rceil\) groups of nearly equal sizes; decreasing the
lower endpoint by one accommodates integer rounding.  Drop interactions
between groups before applying the radial argument.  A position close
to a member of its group is also close to a particle of the original
cube, so the reserve \eqref{en:kinetic-reserve} remains available.
After reducing the harmless lower endpoint to \(m/3\), the group
bound is at least
\[
  \frac{4\pi r}{\ell^3}(1-\delta_D)(m/3-1)k.
\]
Choose once and for all \(b>12\) and then \(C_0>12b\).
Define the convex function on \([0,\infty)\)
\[
 F_b(y)=
 \begin{cases}
 y^2,&0\le y\le b,\\
 2by-b^2,&y>b.
 \end{cases}
\]
Both the small-\(k\) estimate and the group estimate imply the common
bound
\begin{equation}\label{en:smallcube-convex}
 \frac{4\pi r}{\ell^3}(1-\delta_D)
 \left[(D\ell^3)^2F_b\left(\frac{k}{D\ell^3}\right)-k\right].
\end{equation}
For \(k\le m\) this follows from \(F_b(y)\le y^2\).
For \(k>m\), use \(F_b(y)\le2by\) and
\(m/3-1\ge2bD\ell^3-1\) for sufficiently large \(D\).

\emph{From small-cube counts to the desired count penalty.}
Set \(\vartheta=n/(Dv_b)\in[0,6]\).  For all \(y\ge0\),
\begin{equation}\label{en:convex-remainder}
 F_b(y)\ge\vartheta^2+2\vartheta(y-\vartheta)
                         +c\mathfrak h(y-\vartheta).
\end{equation}
If \(y\le b\), the remainder after the tangent is
\((y-\vartheta)^2\).  If \(y>b\), write
\(t=y-\vartheta\) and \(a=b-\vartheta\); the remainder is
\(2at-a^2\ge at\), and \(a\ge b-6>0\).
These two observations prove \eqref{en:convex-remainder} with a
uniform positive constant.  Sum \eqref{en:smallcube-convex} over the
small cubes.  The linear tangent terms cancel because their counts sum
to \(n\).  The subtracted self-count terms cost at most
\[
 C r\ell^{-3}n\le Cv_bD^{0.93}.
\]
The factor \(1-\delta_D\) changes the leading scalar by at most
\(Cv_bD^{1-0.002}\).  Since \(\mathfrak h\) is convex, Jensen's
inequality within each \(d\)-cube gives
\[
 \sum_{\text{small cubes }C\subset A}|C|\,
   \mathfrak h\left(\frac{n_C}{D|C|}-\vartheta\right)
 \ge |A|\,\mathfrak h\left(\frac{n_A}{D|A|}-\vartheta\right).
\]
This proves the scalar and count terms in \eqref{en:rough-ineq},
while leaving \eqref{en:kinetic-reserve} unused.

\emph{Recovery of the one-particle kinetic term.}
We now sum the sector restrictions back on the original \(H^1\) form
domain of the \(E\)-cube; the following estimates concern its global
one-coordinate slices, not an arbitrary function in the larger direct
sum of the cut Neumann form domains.
Let \(u_p\) be the normalized cosine basis of the Neumann cube, with
\(p\in(\pi/E)\Z_{\ge0}^3\).  The vector fields
\(\nabla u_p/|p|\), for \(p\ne0\), are orthonormal.  Let
\(\mathcal P\) be their projection for
\(0<|p|\le D^{0.003}\).  The explicit cosine formulas give
\begin{equation}\label{en:gradient-evaluation}
 \|\mathcal P\|_{L^2\to L^\infty}^2
 \le\frac{C}{E^3}\#\{p:0<|p|\le D^{0.003}\}
 \le CD^{0.009}.
\end{equation}
For a fixed slice in the other particle coordinates, enlarge \(S_i\)
to the union \(S\) of the balls of radius \(s\) about all the other
particles.  Its volume is at most \(Cns^3\).  The integral kernel
bound in \eqref{en:gradient-evaluation}, or duality, implies
\[
 \|\mathcal P\1_S F\|_2^2
       \le CD^{0.009}ns^3\|F\|_2^2
       \le CD^{-0.071+3\theta}\|F\|_2^2.
\]
Thus, for \(F=\nabla f\),
\[
 \|\mathcal P\nabla f\|_2^2
 \le2\|\1_{S^c}\nabla f\|_2^2
                +CD^{-0.071+3\theta}\|\nabla f\|_2^2.
\]
For example, if \(\theta<10^{-3}\), the last coefficient is
\(o(\eps_D)\).  Integration by parts, using the Neumann condition on
\(u_p\), gives
\(\langle\nabla u_p/|p|,\nabla f\rangle
   =|p|\langle u_p,f\rangle\).
A fixed fraction of
\eqref{en:kinetic-reserve} therefore controls
\(\sum_{0<|p|\le D^{0.003}}p^2|\langle u_p,f\rangle|^2\).
Another fraction of the everywhere reserve controls the remaining
weights one, since
\(\eps_Dp^2\ge D^{0.004}\) for \(|p|>D^{0.003}\).
Sum the slice estimates over the coordinates.  Together with the count
bound already obtained this proves \eqref{en:rough-ineq} with
\(c_{\mathrm r}=10^{-3}\).
\end{proof}

\subsection{Localization in the number of low-frequency excitations}
\label{en:localization-section}

The rough estimate now permits a cutoff in an excitation number, in the spirit of the large-matrix localization of \cite[Appendix~A]{LS2001}, without
losing an appreciable energy density.  We record explicitly the form
error because the cutoff need not preserve the hard-core constraint.
It is applied only after passing to the bounded replacement potential.

Let \(P_0=|u_0\rangle\langle u_0|\),
\(n_0=\dd\Gamma(P_0)\), and \(n_+=n-n_0\).  Define
\[
 P_L=\sum_{0<|p|\le K_H/E}|u_p\rangle\langle u_p|,
 \qquad n_L=\dd\Gamma(P_L),\qquad n_H=n_+-n_L,
\]
where \(K_H\asymp Q^{100}\).  Let
\(M\asymp Dv_bQ^{-400}\) be a positive integer.  Choose real
functions \(f,h\) on \(\R\) with
\[
 f^2+h^2=1,\quad f=1\text{ on }(-\infty,M/8],\quad
 f=0\text{ on }[M/4,\infty),\quad
 \operatorname{Lip}(f)+\operatorname{Lip}(h)\le C/M.
\]
For example, write \(f=\cos\vartheta\), \(h=\sin\vartheta\),
where \(\vartheta\) increases smoothly from zero to \(\pi/2\)
on this interval.  Define
\begin{equation}\label{en:three-cutoffs}
 F_1=f(n_L)^2,\qquad F_2=h(n_L)^2,\qquad
 F_3=\sqrt2 f(n_L)h(n_L).
\end{equation}
Then \(\sum_{j=1}^3F_j^2=1\).  The first and third localized
vectors are supported on \(n_L\le M/4\); the second is supported
on \(n_L\ge M/8\).  We also have
\begin{equation}\label{en:count-reconstruction}
 \psi=F_1\psi+F_2\psi,\qquad
 \|A\psi\|^2\le2\|AF_1\psi\|^2+2\|AF_2\psi\|^2
\end{equation}
for every operator \(A\) for which these expressions are defined.
The second inequality is the reason for using three cutoffs: later it
transfers configuration-count estimates without a commutation assumption.

\begin{lemma}[Form error of the excitation cutoff]\label{en:cutoff-error}\label{en:localization}
Let \(\psi\) belong to the form domain of \(H_n\), and let
\(V_n=\sum_{i<j}v(x_i-x_j)\).  Then
\begin{equation}\label{en:cutoff-form-error}
 \left|\sum_{j=1}^3\langle F_j\psi,H_nF_j\psi\rangle
                   -\langle\psi,H_n\psi\rangle\right|
 \le\frac{C}{M^2}\left[
     \langle\psi,V_n\psi\rangle
       +\|v\|_1\frac{K_H^3}{v_b}\,
                  n\langle\psi,n_+\psi\rangle\right].
\end{equation}
For any one-particle operator \(a\) with \(\|a\|\le1\), the
corresponding error for \(\dd\Gamma(a)\) is at most
\(CnM^{-2}\|\psi\|^2\).
\end{lemma}

\begin{proof}
Let \(\Pi_k\) project onto \(n_L=k\), and write
\(\psi_k=\Pi_k\psi\).  The kinetic operator commutes with every
cutoff.  A two-particle multiplication operator changes \(n_L\) by
at most two, so \(\Pi_kV_n\Pi_l=0\) for \(|k-l|>2\).
The difference of the potential forms is exactly
\[
 -\frac12\sum_{k,l}\sum_{j=1}^3
       (F_j(k)-F_j(l))^2\langle\psi_k,V_n\psi_l\rangle.
\]
The cutoffs have Lipschitz constants at most \(C/M\).
Positivity gives
\[
 |\langle\psi_k,V_n\psi_l\rangle|
 \le\langle\psi_k,V_n\psi_k\rangle^{1/2}
      \langle\psi_l,V_n\psi_l\rangle^{1/2}.
\]
Since each row has at most five nonzero blocks, the absolute form
difference is at most
\begin{equation}\label{en:dephased-error}
 \frac{C}{M^2}\langle\psi,\mathcal D(V_n)\psi\rangle,
 \qquad
 \mathcal D(V_n)=\sum_k\Pi_kV_n\Pi_k
       =\frac1{2\pi}\int_0^{2\pi}
          e^{-itn_L}V_ne^{itn_L}\,\dd t.
\end{equation}

The cosine evaluation bound on the low band is
\(\sup_x P_L(x,x)\le CK_H^3/v_b\).  For each pair it implies
\begin{equation}\label{en:pair-low-projection}
 P_{L,i}v(x_i-x_j)P_{L,i}
       \le C\|v\|_1\frac{K_H^3}{v_b}P_{L,i}.
\end{equation}
Indeed, condition on \(x_j\) and all other coordinates, apply
\(|P_Lg(x_i)|^2\le(\sup_xP_L(x,x))\|P_Lg\|_2^2\), and
integrate the potential in \(x_i\).
For the pair \((i,j)\), the rotations on the other coordinates cancel.
Writing \(z=e^{it}-1\), the remaining rotation is
\[
 (1+zP_{L,i})(1+zP_{L,j})
   =1+zP_{L,i}+zP_{L,j}+z^2P_{L,i}P_{L,j}.
\]
The squared norm after multiplication by \(v_{ij}^{1/2}\) is bounded
by a fixed constant times the sum of the four squared norms.
Use \eqref{en:pair-low-projection} on the last three terms, and
\(P_{L,i}P_{L,j}\le P_{L,i}\), to obtain
\[
 \langle\psi,\mathcal D(V_n)\psi\rangle
 \le C\langle\psi,V_n\psi\rangle
      +C\|v\|_1\frac{K_H^3}{v_b}\,
                       n\langle\psi,n_L\psi\rangle.
\]
This and \(n_L\le n_+\) prove
\eqref{en:cutoff-form-error}.  A one-particle operator has block
width one instead of two.  Using its second-quantized norm at most
\(n\), and \(2\|\psi_k\|\|\psi_l\|\le
\|\psi_k\|^2+\|\psi_l\|^2\), proves the final assertion.
\end{proof}

We apply the lemma to the restrictions of a normalized hard-core ground
vector to box-assignment sectors.  Its zero extension belongs to the
soft-potential form domain, and the potential expectation is bounded by
the hard-core energy.  The localized vectors need not satisfy the hard-core
constraint; all their lower bounds use \(H_n\).  The finite-rank
cosine projections preserve \(H^1\), so the localized vectors are
in this form domain.

Here is the quantitative energy accounting.  Write
\(e_0(x)=4\pi rx^2\).  Summing the rough bound and comparing with the
thermodynamic upper bound
\(V(e_0(D)+O(1))\) gives, over boxes with \(n\le6Dv_b\),
\begin{equation}\label{en:preliminary-excitations}
 \frac1V\sum\langle n_+\rangle
           \le CQ^2D^{1-c_{\mathrm r}}.
\end{equation}
To justify using the tangent to \(e_0\) in this summation for every
particle sector, note first that \(n\le6Dv_b\) gives the exact scalar
remainder \(4\pi r(n-Dv_b)^2/v_b\).  If \(n>6Dv_b\), partition
the labels into groups of sizes between \(3Dv_b-o(Dv_b)\) and
\(6Dv_b\), apply Lemma~\ref{en:rough-bound} to each group and drop
intergroup interactions.  The energy per particle is then at least
\(12\pi\alpha-o(1)\), whereas the slope of the tangent is
\(8\pi\alpha\).  These sectors have a positive excess of order
\(n\) above that tangent.  Finally,
\(\min\{p^2,1\}\ge cE^{-2}\) for \(p\ne0\), which proves
\eqref{en:preliminary-excitations}.  All sums include the squared norms
of sector vectors and any spectator variables.

We use the cutoffs only for \(Dv_b/2\le n\le6Dv_b\).
The total potential expectation is at most \(CDV\), and
\(M\asymp DQ^{-397}\).  Substitution of
\eqref{en:preliminary-excitations} into
\eqref{en:cutoff-form-error} gives the explicit bound
\begin{equation}\label{en:localization-density-error}
 \frac1V\sum\left|\text{localization form error}\right|
 \le C\left(D^{-1}Q^{794}+D^{-c_{\mathrm r}}Q^{1097}\right)
 \le CQ^{-1},
\end{equation}
provided \(795\theta<1\) and \(1098\theta<c_{\mathrm r}\).
The one-particle observable error is bounded by the first term on the
right.  Thus these errors are negligible on the energy scale used below.



\subsection{The sharp bound on a low-excitation piece}

The next bound is the information beyond the scalar LHY formula that
will determine the mesoscopic occupation.  We give the passage from the
published operator estimates to the exact vectors, including the
zero-mode resolution and the cosine-basis convention.

For $z\ne0$ put $\zeta_z=z/|z|$, and set $\zeta_0=1$.  Define
\begin{align}
 \tau(p)&=p^2-\frac\pi{2E^2}-\frac{K_H}{E^2}
                       \1_{\{p\notin\mathcal P_L\}},\label{en:tau}\\
 \rho_z&=|z|^2/v_b,\qquad
 D_p(z)=\sqrt{\tau(p)^2+2\rho_z\widehat g(p)\tau(p)},\nonumber\\
 \beta_p(z)&=\frac{\tau(p)+\rho_z\widehat g(p)-D_p(z)}
                   {\rho_z\widehat g(p)},\qquad
 b_p(z)=\frac{a_p+\zeta_z^2\beta_p(z)a_p^*}{\sqrt{1-\beta_p(z)^2}}.
 \label{en:bogoliubov}
\end{align}
At zero numerator and denominator use the continuous value.  The factor
$\zeta_z^2$ restores the phase in the pairing term; for real $z$ it
equals one.  Put $\widetilde D_p=D_p$ on $\mathcal P_L$ and
$\widetilde D_p=K_H^{-1}D_p$ on its complement.  The square root is real:
the support of $g$ gives
$|\widehat g(p)-\widehat g(0)|\le Cr^3p^2$, and
$\tau(p)\asymp p^2$.  In particular the possible negative part of
$\widehat g$ cannot destroy positivity when $D$ is large and
$|z|^2\le2n$ in \eqref{en:middle-sector}.
More explicitly, $\rho_z\le CD$ implies
$\tau(p)+2\rho_z\widehat g(p)\ge(c-CD^{-2})p^2$.
At $\rho_z=0$ we have $D_p=\tau(p)$, $\beta_p=0$, and $b_p=a_p$.

For $n$ in \eqref{en:middle-sector} and an $n$-particle vector
$\psi$, let
$\psi_z=(\langle z|\otimes1)\psi$ denote its zero-mode coherent
projection, with measure $\dd^2z/\pi$.  Here
$|z\rangle=e^{-|z|^2/2}\sum_{k\ge0}z^k|k\rangle/\sqrt{k!}$.
Define the nonnegative quadratic expression
\begin{equation}\label{en:bog-energy}
 \mathcal B(\psi)=\int_{|z|^2\le2n}
       \sum_{p\ne0}\widetilde D_p(z)\|b_p(z)\psi_z\|^2
                                  \frac{\dd^2z}{\pi}.
\end{equation}

For later truncations we record the coherent tail at this point.
There is no large-$z$ assumption on the original vector.  Write
$\psi=\sum_{k=0}^n|k\rangle\otimes\psi^{(n-k)}$.  Angular integration
shows that its radial coherent density is a sum of gamma densities
$e^{-s}s^k/k!$, $0\le k\le n$, weighted by
$\|\psi^{(n-k)}\|^2$.  The gamma Chernoff bound consequently gives
\begin{equation}\label{en:coherent-tail}
 \int_{|z|^2>2n}(1+|z|)^A\|\psi_z\|^2
       \frac{\dd^2z}{\pi}\le C_A n^{A/2}e^{-cn}\|\psi\|^2
\end{equation}
for every fixed $A\ge0$.
The local lower bound we shall prove is the following.

\begin{lemma}[Energy with a retained quasiparticle term]\label{en:sharp}
For $n$ in \eqref{en:middle-sector} and a bosonic $n$-particle
vector $\psi$ in the form domain of $H_b$ satisfying
$\psi=\1_{\{n_L\le M/4\}}\psi$,
\begin{equation}\label{en:sharp-bound}
 \langle\psi,H_b\psi\rangle
 \ge\{v_be_r(n/v_b)-Cv_bQ^{-1/4}\}\|\psi\|^2
       +cE^{-2}\langle n_++K_Hn_H\rangle_\psi
       +c\mathcal B(\psi).
\end{equation}
The constants are uniform in the particle sector and in spectator
variables.
\end{lemma}

Its proof must preserve two positive terms: the gap controlling physical
excitation numbers, and the coherent quadratic energy $\mathcal B$.
Write $S=E^{-2}(\pi n_+/2+K_Hn_H)$.  The symmetrization and
second-quantization formulas of \cite[Theorem~2.6 and Lemma~2.8]{FJGMOT}
separate the number-operator block
$\widehat g(0)\{n(n-1)-n_+(n_+-1)\}/(2v_b)$ from terms containing
the constant mode.  Coherent resolution of the latter terms gives the
quadratic--cubic operator $\mathcal H_0(z)$ defined below, with the
upper-symbol corrections accounted for in the proof of
Lemma~\ref{en:sharp}.  We first bound this operator while retaining
its excitation energy, then estimate the correction to the cosine
coefficients, and finally return to the box Hamiltonian.

The next operator estimate specifies the gap fractions in the energy
analysis of \cite{FJGMOT}.
Write $\mathcal P_L^{\mathbb Z}$ for the signed version of the low band,
extend operator subscripts by componentwise absolute value, put
$c_p=\prod_i c_{p_i}$, and let
$c(p,k)=\prod_i c_{k_i-p_i}/(c_{p_i}c_{k_i})$ be the printed coefficient
in that paper.  Define
\begin{align}
 Q_3(z)&=\frac1{v_b}\sum_{\substack{p\in\mathcal P_L^{\mathbb Z},\ k\ne0\\
                                      |p-k|\ne0}}
 c(p,k)\widehat g(k)(\overline z a_p^*a_{|p-k|}a_k+\mathrm{h.c.}),
                                                        \label{en:printed-cubic}\\
 \mathcal H_0(z)&=\sum_{p\ne0}\left[
   (\tau(p)+\rho_z\widehat g(p))a_p^*a_p
   +\frac{\widehat g(p)}{2v_b}
      (\overline z^2a_p^2+z^2a_p^{*2})\right]
   +\frac{\widehat{g\omega}(0)|z|^4}{2v_b}\nonumber\\
 &\quad+\rho_z\sum_{p\ne0}
       \{\widehat{g\omega}(0)+\widehat{g\omega}(p)\}a_p^*a_p+Q_3(z).
                                                        \label{en:H0}
\end{align}
Here $k$ ranges over the nonnegative Neumann lattice, and $|p-k|$ in
an operator subscript means componentwise absolute value.
Thus $\mathcal H_0$ is (2.36) of \cite{FJGMOT} at zero chemical
parameter with its artificial scalar convexity term subtracted.
The exact cosine coefficients are treated separately in
Lemma~\ref{en:cosine-correction}.

\begin{lemma}[A small-gap version of the quadratic--cubic bound]
\label{en:small-gap}
Let $n$ satisfy \eqref{en:middle-sector}, $|z|^2\le2n$, and
$\xi=\1_{\{n_+\le n,\ n_L\le M/4\}}\xi$.  For every fixed $\delta>0$
and sufficiently large $D$,
\begin{align}
 \langle\xi,\mathcal H_0(z)\xi\rangle
 &\ge 4\pi c_{\rm L}(a')^{5/2}v_b\rho_z^{5/2}\|\xi\|^2
       +\sum_{p\ne0}\widetilde D_p(z)\|b_p(z)\xi\|^2
       -\delta\langle S\rangle_\xi
       -C_\delta v_bQ^{-1/4}\|\xi\|^2 .
                                                        \label{en:small-gap-bound}
\end{align}
\end{lemma}

\begin{proof}
We use the Bogoliubov identity (2.39), the completion of a square
(7.18)--(7.20), and the scalar integral estimate in
\cite[Lemma~A.1]{FJGMOT}.  We give the intervening bounds, in particular
an estimate of the creation term with every commutator retained.
This obtains the required gap fraction directly; no quantifier about
the coefficient in that paper's free-energy Theorem~2.9 is needed.
By homogeneity suppose $\|\xi\|=1$, and abbreviate
$K=K_H$, $N_+=\langle n_+\rangle_\xi$,
$\lambda=(Dv_b)^{-1/2}$, and
$h(k)=|\widehat g(k)|/(8\pi a')\le1$.

The Bogoliubov identity (2.39) of \cite{FJGMOT} gives the exact
separation
\begin{equation}\label{en:H0-decomposition}
 \mathcal H_0(z)=\mathcal E_{\rm B}(z)
       +\sum_{p\ne0}D_p(z)b_p(z)^*b_p(z)
       +\mathcal R_{\rm int}(z)+Q_3(z),
\end{equation}
where
\begin{align*}
 \mathcal E_{\rm B}(z)
   &=\frac12\sum_{p\ne0}
          \bigl(D_p(z)-\tau(p)-\rho_z\widehat g(p)\bigr)
            +\frac{\widehat{g\omega}(0)|z|^4}{2v_b},\\
 \mathcal R_{\rm int}(z)
   &=\rho_z\sum_{p\ne0}
       \{\widehat{g\omega}(0)+\widehat{g\omega}(p)\}a_p^*a_p.
\end{align*}
We shall bound $Q_3$ using part of the positive quadratic energy,
cancel its normal-ordering remainder against $\mathcal R_{\rm int}$,
and evaluate the scalar $\mathcal E_{\rm B}$ at the end.

First remove the part of \eqref{en:printed-cubic} with $k$ low.
Put $l=|p-k|$.  Bounded multiplicity of $k\mapsto l$ for fixed $p$,
commutation of $a_l$ past $a_l^*$, and $n_L\le M/4$ give
\[
 \sum_{p,k}\|a_l^*a_p\xi\|^2\le C(M+K^3)N_+,
 \qquad
 \sum_{p,k}\|a_k\xi\|^2\le CK^3N_+.
\]
Here both $p$ and $k$ range over their low bands, with the signed
multiplicity in $p$ bounded by eight; the terms with $l=0$ are omitted.
Cauchy--Schwarz in the expectation
$\langle a_l^*a_p\xi,a_k\xi\rangle$ therefore bounds this part by
\[
 C\lambda K^{3/2}(M+CK^3)^{1/2}N_+.
\]
Its coefficient relative to $E^{-2}N_+$ is
$O(Q^{-48}+D^{-1/2}Q^{300.5})=o(1)$.

In the remaining cubic term $k$ is high.  Substitute
$a_k=U_kb_k+V_kb_k^*$ and write it as $\mathcal T_1+\mathcal T_V$.
For high $k$, $|U_k|\le C$ and
$|V_k|\le Ch(k)/k^2$.  Thus the coefficients of $\mathcal T_V$
are at most $C\lambda h(k)^2/k^2$ in absolute value.
Let $l=|p-k|$.  Since $p$ is low and $k$ high, the exact commutation is
\begin{equation}\label{en:Talpha-commute}
 a_p^*a_l b_k^*=b_k^*a_p^*a_l
                         +U_k\1_{\{l=k\}}a_p^* .
\end{equation}
The harmless common phase of $z$ can be included in the coefficients.
Set
\[
 B_H=\sum_{k\ {\rm high}}\frac{k^2}{K}\|b_k\xi\|^2.
\]
For the first term on the right of \eqref{en:Talpha-commute},
Cauchy--Schwarz in $p$, followed by weighted Cauchy--Schwarz in $k$,
gives
\begin{align}
 &C\lambda\sum_{\substack{k\ {\rm high}\\p\ {\rm low,\ signed}}}
       \frac{h(k)^2}{k^2}\|b_k\xi\|\,\|a_p^*a_l\xi\|\nonumber\\
 &\qquad\le
 C\lambda K^2 B_H^{1/2}
  \left(\sum_{p,k}\frac{h(k)^4}{k^6}
                  \|a_p^*a_l\xi\|^2\right)^{1/2}.
                                                        \label{en:Talpha-CS}
\end{align}
The summation ranges on the last line are those on the first.
The identity
$\|a_p^*a_l\xi\|^2=\|a_pa_l\xi\|^2+\|a_l\xi\|^2$
and bounded folding multiplicity imply
\[
 \sum_{p,k}\|a_pa_l\xi\|^2\le CMN_+,\qquad
 \sum_{p,k}\|a_l\xi\|^2\le CK^3N_+.
\]
For the first bound use
$\sum_p a_p^*n_+a_p\le Cn_Ln_+$; it holds also when $l=p$.
Since $k^2\ge K^2/E^2$, Young's inequality in
\eqref{en:Talpha-CS} yields, for every $\varepsilon>0$,
\begin{align}
 |\langle\mathcal T_V^{\,\rm reordered}\rangle|
 &\le\varepsilon B_H+
 C_\varepsilon\lambda^2K^4(E/K)^6(M+K^3)N_+\nonumber\\
 &\le\varepsilon B_H+
 C_\varepsilon(Q^{-594}+D^{-1}Q^{103})N_+.
                                                        \label{en:Talpha-main-bound}
\end{align}

The second term in \eqref{en:Talpha-commute} is also summable.
The condition $|p-k|=k$ requires, in each coordinate,
$p_i=0$ or $p_i=2k_i$.  Some $p_i$ is nonzero, so for fixed $p$
the admissible $k$ lie in at most three planes $k_i=p_i/2$.
Two-dimensional lattice summation, retaining $h(k)^2$, gives
\[
 \sum_{\substack{k\ {\rm high}\\|p-k|=k}}\frac{h(k)^2}{k^2}
       \le CE^2(1+\log D).
\]
Indeed the dyadic sums below a sufficiently large polynomial frequency
are $O(E^2)$ per annulus, and the radial transform bound
$h(k)\le D^C/|k|$ makes the remaining tail summable.
Cauchy--Schwarz in the $O(K^3)$ possible values of $p$ now gives
\begin{align}
 |\langle\mathcal T_V^{\,\rm commutator}\rangle|
 &\le C\lambda E^2K^{3/2}(1+\log D)\sqrt{N_+}\nonumber\\
 &\le\delta\langle S\rangle+
                   C_\delta D^{-1}Q^{303}(1+\log D)^2.
                                                        \label{en:Talpha-resonance}
\end{align}
This ordering of the operators avoids an unweighted sum of vacuum
commutators over all high modes.  Choose $\varepsilon$ small enough
that $\varepsilon B_H\le K^{-1}\sum_{k\ {\rm high}}D_k\|b_k\xi\|^2$.
The other terms in \eqref{en:Talpha-main-bound} are $o(\langle S\rangle)$;
the scalar in \eqref{en:Talpha-resonance} is $o(v_bQ^{-1})$.
The coefficient $\delta$ here can be any prescribed fixed fraction of
the one in the statement.

It remains to treat $\mathcal T_1$.  Complete a square with
$(1-2K^{-1})\sum_{k\ {\rm high}}D_kb_k^*b_k$.
The square is nonnegative.  Its negative remainder, after normal
ordering, is the sum of a commutator term $\mathcal T_c$ and a quartic
term $\mathcal T_0$, as in (7.18)--(7.20) of \cite{FJGMOT}.
The coefficient in that square is
\[
 A_k=\frac{z\widehat g(k)}
 {v_b(1-2K^{-1})D_k\sqrt{1-\beta_k^2}}
       \sum_{\substack{p\in\mathcal P_L^{\mathbb Z}\\|p-k|\ne0}}
                  c(p,k)a_{|p-k|}^*a_p.
\]
In $A_k^*A_k$, the commutator is nonzero only when
$|p-k|=|s-k|$.  Either $p=s$, or at least one coordinate of $k$
lies in the low-frequency slab $|k_i|\le K/E$.  The exact signed-cosine
weight identity $\sum_p c_p^{-2}a_p^*a_p=n_L$, followed by
\cite[Lemmas~7.2--7.3, Equation~(7.27)]{FJGMOT}, gives
\begin{equation}\label{en:Tc-bound}
 \langle\mathcal T_c\rangle
 \le2\rho_z\widehat{g\omega}(0)N_+
       +C(Q^{-97.75}+D^{-1}Q^{200.25})E^{-2}N_+ .
\end{equation}
The first term cancels the remaining quadratic interaction term in
\eqref{en:H0}, up to an error.  Splitting that interaction into its
low and high frequencies, Lemma~7.2(1) of the same paper bounds this
error by
\[
 Cn_H+CD^{-2}Q^{198}n_+
       \le C(Q^{-98}+D^{-2}Q^{200})S.
\]
The quartic term has the explicit bound
\[
 \mathcal T_0\le C\,\frac{a'K_\ell^3K}{v_b}n_Ln_+,
\]
obtained by Cauchy--Schwarz in its two low indices and summing the
remaining high occupation, as in (7.28)--(7.29) of \cite{FJGMOT}.
On the present subspace its ratio to $E^{-2}n_+$ is at most
$CMa'K_\ell^3K/E=O(Q^{-295})$.
All errors in this paragraph are therefore $o(S)$.
These estimates are uniform for $0\le\rho_z\le2n/v_b$.
Lemma~7.2 of \cite{FJGMOT} concerns the potential alone, and
Lemma~7.3 is applied at $\rho_{\rm aux}$ and multiplied by
$\rho_z/\rho_{\rm aux}$, which is uniformly bounded.
At $\rho_z=0$ the cubic and interaction terms vanish and $b_p=a_p$.

In the notation of \eqref{en:H0-decomposition}, the cancellation
just proved reads
\[
 \langle\mathcal R_{\rm int}-\mathcal T_c\rangle
       \ge-o(1)\langle S\rangle.
\]
The completed square spends $(1-2K^{-1})$ of the high quadratic
energy; the reordered creation term spends at most another $K^{-1}$.
Together with the low-cubic and quartic bounds this leaves
\[
 \langle\mathcal H_0(z)\rangle
   \ge\mathcal E_{\rm B}(z)
       +\sum_{p\ne0}\widetilde D_p(z)\|b_p(z)\xi\|^2
       -\delta\langle S\rangle-C_\delta v_bQ^{-1}.
\]
We have retained every low Bogoliubov mode and at least a fraction
$K^{-1}$ of every high mode.  Lemma~A.1 of \cite{FJGMOT} evaluates
$\mathcal E_{\rm B}(z)$ as
$4\pi c_{\rm L}(a')^{5/2}v_b\rho_z^{5/2}$ with error
$Cv_bQ^{-1/4}$, uniformly on $|z|^2\le2n$.  Combining the two bounds,
choosing the fixed gap fractions first and then increasing $D$, proves
\eqref{en:small-gap-bound}.
\end{proof}

\begin{lemma}[Exceptional cosine coefficients]\label{en:cosine-correction}
Let $n$ satisfy \eqref{en:middle-sector}, let $|z|^2\le2n$, and restrict
the excitation space to $n_+\le n$ and $n_L\le M/4$.  On this space the
exact cosine expansion may replace the coefficients in (2.35)--(2.37)
of \cite{FJGMOT}.  Its correction is bounded in absolute quadratic form by
\[
 \delta\left(S+\sum_{p\ne0}\widetilde D_p(z)b_p(z)^*b_p(z)\right)
       +C_\delta v_bQ^{-1/4}
\]
for every fixed $\delta>0$, if $D$ is sufficiently large.
\end{lemma}

\begin{proof}
Extend a cosine subscript by componentwise absolute value.  In a term
$a_p^*a_l a_k$, $p$ is a signed low frequency and $l=|p-k|$; terms with
$l=0$ are omitted.  The product-to-sum identity in each coordinate shows
that the printed coefficient agrees with the exact expansion unless
$k_i=0$ or $k_i=|p_i|$ for some $i$.  All multiplicities and discrepancies
are bounded by an absolute constant.  For example, when $k_i=0$ and
$p_i\ne0$, the two choices of the sign of $p_i$ must share one coefficient;
when $k_i=p_i\ne0$ and $l_i=0$, the zero cosine has normalization one.
These are precisely the two exceptional cases.

Put $\lambda=(Dv_b)^{-1/2}$ and
$h(k)=|\widehat g(k)|/(8\pi a')\le1$.  Every correction has coefficient
bounded by $C\lambda h(k)$.  We estimate its expectation in a normalized
excitation vector $\xi$ with $n_+\le n$, $n_L\le M/4$.
For $|k|\le3K_H/E$, Cauchy--Schwarz applied to
$\langle a_l^*a_p\xi,a_k\xi\rangle$ gives
\begin{equation}\label{en:cosine-low}
 C\lambda K_H^{3/2}
       \{M\langle n_+\rangle+CK_H^3\langle n_+\rangle\}^{1/2}
       \langle n_+\rangle^{1/2}.
\end{equation}
To see the first factor, sum the occupation products after commuting
$a_l$ past $a_l^*$; for each $p$, the multiplicities of $l$ are bounded.
The second Cauchy--Schwarz factor sums $a_k^*a_k$ at most $CK_H^3$
times.  The two coefficients multiplying $\langle n_+\rangle$ in
\eqref{en:cosine-low} are $O(Q^{-50})$ and
$O(D^{-1/2}Q^{298.5})$, respectively, and are $o(E^{-2})$.

For $|k|>3K_H/E$, both $k$ and $l$ are high frequencies.  Use
$a_k=U_kb_k+V_kb_k^*$, where $|U_k|\le C$ and $|V_k|\le C/k^2$.
The terms involving $\|b_k\xi\|$, including the excitation part of
$\|b_k^*\xi\|$, are bounded by the product of
\[
 \left(\sum_k k^2K_H^{-1}\|b_k\xi\|^2\right)^{1/2}
\]
and a factor whose square is at most
\begin{equation}\label{en:cosine-high}
 C\lambda^2K_H^4\left\{
 (E/K_H)^2M\langle n_H\rangle
       +E^2(1+\log D)\langle n_+\rangle\right\}.
\end{equation}
Here is the complete weighted sum underlying this estimate.  Restrict
$(p,k)$ to the exceptional planes of the correction and put
$F_{pk}=\|a_l^*a_p\xi\|$, with $l=|p-k|$.
Cauchy--Schwarz in the at most $CK_H^3$ values of $p$, then in $k$
with weight $k^2/K_H$, gives
\[
 C\lambda\sum_k h(k)\|b_k\xi\|\sum_pF_{pk}
 \le C\lambda K_H^2
 \left(\sum_k\frac{k^2}{K_H}\|b_k\xi\|^2\right)^{1/2}
 \left(\sum_{p,k}\frac{h(k)^2}{k^2}F_{pk}^2\right)^{1/2}.
\]
Since $l$ is high and $p$ low,
$F_{pk}^2=\|a_la_p\xi\|^2+\|a_p\xi\|^2$ and bounded folding
multiplicity gives
\begin{align*}
 \sum_{p,k}\frac{h(k)^2}{k^2}\|a_la_p\xi\|^2
       &\le C(E/K_H)^2\langle n_Ln_H\rangle
        \le C(E/K_H)^2M\langle n_H\rangle,\\
 \sum_{p,k}\frac{h(k)^2}{k^2}\|a_p\xi\|^2
       &\le CE^2(1+\log D)\langle n_+\rangle.
\end{align*}
The second line uses the two-dimensional plane sums
\begin{equation}\label{en:plane-sums}
 \sum_{\substack{k\text{ on an exceptional plane}\\|k|>3K_H/E}}
       \frac{h(k)^i}{k^2}\le CE^2(1+\log D),\qquad i=1,2.
\end{equation}
For \eqref{en:plane-sums}, use dyadic annuli up to a sufficiently large
power of $D$.  Beyond that scale the radial transform formula
$\widehat g(k)=4\pi|k|^{-1}\int_0^r t g(t)\sin(|k|t)\dd t$
gives $h(k)\le D^C/|k|$ and a summable tail.  The same argument with
$k^{-4}$ gives a bound $CE^2$.

The extra unit in $\|b_k^*\xi\|^2=\|b_k\xi\|^2+1$ comes with $V_k$.
Its total contribution is at most
\begin{equation}\label{en:cosine-creation}
 C\lambda K_H^{3/2}E^2
       \left\{\sqrt{M\langle n_H\rangle}
                 +(1+\log D)\sqrt{\langle n_+\rangle}\right\}.
\end{equation}
To verify \eqref{en:cosine-creation}, bound the unit contribution by
$C\lambda\sum_{p,k}h(k)k^{-2}F_{pk}$.
Its occupation-product part is at most
\[
 C\lambda
 \left(\sum_{p,k}\frac{h(k)^2}{k^4}\right)^{1/2}
 \left(\sum_{p,k}\|a_la_p\xi\|^2\right)^{1/2}
 \le C\lambda E^2K_H^{3/2}\sqrt{M\langle n_H\rangle}.
\]
Here we used the loose bound $CE^4$ for each $k^{-4}$ plane sum
and at most $CK_H^3$ values of $p$.
Its commutator part is bounded by first summing
$h(k)/k^2$ on the plane and then using
$\sum_p\|a_p\xi\|\le CK_H^{3/2}\sqrt{\langle n_+\rangle}$.
These are the two terms in \eqref{en:cosine-creation}.

Young's inequality absorbs the first factor above into a $\delta$
fraction of the high-frequency quasiparticle energy.  In
\eqref{en:cosine-high}, the coefficients of $\langle n_H\rangle$ and
$\langle n_+\rangle$ are $O(Q^{-198})$ and
$O(D^{-1}Q^{399}(1+\log D))$.  They are negligible relative to the gaps
$Q^{98}n_H$ and $Q^{-2}n_+$.  Applying Young's inequality to
\eqref{en:cosine-creation} leaves scalar errors bounded by
$C_\delta[Q^{-194}+D^{-1}Q^{303}(1+\log D)^2]$.
They fit within $C_\delta v_bQ^{-1/4}$ for sufficiently small $u$.
This proves the stated bound.  On the omitted large-$z$ region the
correction has a polynomial bound on the truncated particle sectors:
the $k$ coefficients have a polynomially bounded lattice square sum,
and $p$ lies in a finite band.  Thus \eqref{en:coherent-tail} also covers
that region.
\end{proof}



\begin{proof}[Proof of Lemma~\ref{en:sharp}]
We detail both the gap allocation and the resolution of the zero mode.
The calculations may first be made with finitely many Neumann modes.
These truncations commute with the excitation cutoffs and approximate
the local form domain.  The bounded-potential terms then converge in
form, and the retained positive quasiparticle expression passes by lower
semicontinuity.  This justifies the infinite sums below.
Set
\[
 S=E^{-2}(\pi n_+/2+K_Hn_H),\qquad H_b=H_b^{\rm mod}+S.
\]
Apply the symmetrization estimate of \cite[Theorem~2.6]{FJGMOT}
with the fixed choice $\varepsilon=1/10$.  It gives
\begin{equation}\label{en:symmetrization}
 H_b^{\rm mod}\ge H_b^{\rm sym}-CnDr(r/E)-\varepsilon G,
 \qquad
 G=\tfrac12(1+n_L/M)S.
\end{equation}
On the present vector $G\le5S/8$, so $\varepsilon G\le S/16$.
The hypotheses with the constant associated to this fixed choice hold
for large $D$, since $K_H/K_\ell^4\asymp Q^{96}$ and
$K_\ell K_H^3\sqrt{\rho_{\rm aux}(a')^3}=O(D^{-1}Q^{301})$.
The scalar error is $O(E^2)$ and hence fits within $Cv_bQ^{-1/4}$.

In the second-quantized identity of \cite[Lemma~2.8]{FJGMOT}, the terms
involving only the number operators are
\[
 \frac{\widehat g(0)}{2v_b}
       \{n(n-1)-n_+(n_+-1)\}.
\]
The negative term can be paid from an arbitrarily small fraction of $S$.
Indeed
\begin{equation}\label{en:self-interaction}
 n_+^2\le2Mn_++2nn_H,
 \qquad
 \frac r{v_b}n_+^2
   \le C\{Q^{-400}n_++n_H\}=o(1)S.
\end{equation}
The replacement of $n(n-1)$ by $n^2$ costs $O(1)$.

The coherent resolution is an exact identity for quadratic forms; this
use of upper symbols is the one in \cite{LSY2005}.
The upper symbols of the two zero-mode monomials are
\begin{equation}\label{en:upper-symbols}
 a_0^*a_0:\ |z|^2-1,\qquad
 a_0^{*2}a_0^2:\ |z|^4-4|z|^2+2.
\end{equation}
Monomials containing only creations or only annihilations have their
ordinary polynomial symbols.  Let $\mathcal C(z)$ denote the exact
cosine cubic operator minus the printed cubic operator
\eqref{en:printed-cubic}; Lemma~\ref{en:cosine-correction} bounds this
difference.  If $x=|z|^2$, the upper-symbol correction is
\begin{equation}\label{en:symbol-error}
 \mathcal R(z)=\frac{\widehat{g\omega}(0)}{2v_b}(-4x+2)
 -\frac1{v_b}\sum_{p\ne0}
  \{\widehat g(p)+\widehat{g\omega}(0)+\widehat{g\omega}(p)\}a_p^*a_p.
\end{equation}
The complete form of the second-quantized identity
\cite[Lemma~2.8, Equation~(2.33)]{FJGMOT}, with the exact cosine
coefficients, is therefore
\begin{align}
 \langle\psi,H_b^{\rm sym}\psi\rangle
 &=\frac{\widehat g(0)}{2v_b}
       \left[n(n-1)\|\psi\|^2
                     -\langle\psi,n_+(n_+-1)\psi\rangle\right]
       \nonumber\\
 &\quad+\int_{\mathbb C}
       \langle\psi_z,
          (\mathcal H_0(z)+\mathcal R(z)+\mathcal C(z))\psi_z\rangle
                                      \frac{\dd^2z}{\pi}.
 \label{en:coherent-Hsym-identity}
\end{align}
Here $\mathcal H_0$ contains the modified kinetic symbol $\tau$ and
the low-band cubic cutoff from \eqref{en:H0}; the artificial scalar
convexity term in (2.36) of that paper is absent.  The identity has no
additional remainder: the interaction pieces removed before
symmetrization are already included in the error of
\eqref{en:symmetrization}.
The integrated absolute quadratic form of $\mathcal R$ is bounded by
\[
 C\frac r{v_b}\int\left[
       (1+|z|^2)\|\psi_z\|^2+\langle\psi_z,n_+\psi_z\rangle\right]
                                    \frac{\dd^2z}{\pi}
 \le C\|\psi\|^2.
\]
The last inequality follows from
$\int|z|^2\|\psi_z\|^2\dd^2z/\pi
=\langle n-n_++1\rangle_\psi$ and $rn/v_b=O(1)$.

Apply Lemma~\ref{en:small-gap} to the operator
$\mathcal H_0(z)$ on $|z|^2\le2n$, with a fixed sufficiently small
$\delta$.  The coherent projection preserves both excitation cutoffs.
The lemma retains $\sum\widetilde D_pb_p^*b_p$ and the LHY scalar at
$\rho_z$, while using at most $\delta S$.  Together with
$G\le5S/8$ in \eqref{en:symmetrization} and the $o(S)$ loss in
\eqref{en:self-interaction}, this leaves a fixed positive fraction of
$S$.  In terms of the published formula (2.36), the artificial scalar
convexity term has been subtracted identically, not estimated by its
sign.  Our argument uses no trace inequality from that paper's
Theorem~2.9.

By \eqref{en:coherent-tail}, the discarded coherent region has
exponentially small weighted mass.  All negative terms of the upper symbol have
polynomial norm bounds in $D,n,z$ on excitation sectors at most $n$;
this follows either from the bounded potential or from its projected
second-quantized blocks.  Dropping positive kinetic energy, the omitted
region thus has exponentially small lower-bound error.

It remains to compare the LHY scalar at $\rho_z$ with that at $n/v_b$.
Convexity and its derivative, bounded by $C/D$ in the relevant density
range, give an integrated loss at most
$CD^{-1}(1+\langle n_+\rangle_\psi)$, apart from the negligible tail.
This too is absorbed by the retained gap.  Replacing $a'$ by $r$ costs
$O(v_bQ^{-1})$ by \eqref{en:replacement}.

The cosine convention in the printed cubic formula requires one final
check.  Lemma~\ref{en:cosine-correction} shows that using exact
cosine products costs only a prescribed small fraction of $S$ and of
the quasiparticle term, together with $Cv_bQ^{-1/4}$.  Combining these
estimates leaves fixed positive fractions of both terms and proves
\eqref{en:sharp-bound}.
\end{proof}

